\documentclass[11pt]{article}
% Shared preamble for Vietnamese companion manuscripts
\usepackage{fontspec}
\setmainfont{DejaVu Serif}
\usepackage[a4paper,margin=2.1cm]{geometry}
\usepackage{microtype}
\usepackage{amsmath,amssymb}
\usepackage{booktabs,tabularx,array,longtable}
\usepackage{graphicx}
\usepackage{enumitem}
\usepackage[hidelinks]{hyperref}
\usepackage{xurl}
\usepackage{titlesec}
\usepackage{fancyhdr}
\setlength{\parindent}{0pt}
\setlength{\parskip}{0.55em}
\setlist[itemize]{leftmargin=1.5em,itemsep=0.25em}
\setlist[enumerate]{leftmargin=1.7em,itemsep=0.25em}
\renewcommand{\arraystretch}{1.12}
\titleformat{\section}{\large\bfseries}{\thesection}{0.6em}{}
\titleformat{\subsection}{\normalsize\bfseries}{\thesubsection}{0.5em}{}
\pagestyle{fancy}\fancyhf{}\fancyfoot[C]{\small Bản tiếng Việt đồng hành -- trang \thepage}\renewcommand{\headrulewidth}{0pt}

\emergencystretch=3em
\sloppy

\title{\textbf{GovMut-Health: Đánh giá khả năng phát hiện lỗi quản trị phụ thuộc chuỗi thao tác bằng kiểm thử đột biến}}
\author{Nguyễn Ngọc Thiện, Trịnh Minh Quang \\ Trường THPT Hai Bà Trưng, Huế, Việt Nam}
\date{2026}
\begin{document}\maketitle
\begin{abstract}
Nhiều lỗi quản trị trong AI y tế có trạng thái chỉ xuất hiện sau một chuỗi thao tác, thay vì ở một lời gọi đơn lẻ. Chẳng hạn, hệ thống có thể xử lý đúng khi cung cấp dữ liệu cho mô hình, đúng khi kiểm tra đồng thuận và đúng khi tạo đề xuất, nhưng vẫn chấp nhận một thao tác ghi sai nếu đồng thuận đã thay đổi, đề xuất cũ được thử lại hoặc dữ liệu trong bộ nhớ đệm bị tái sử dụng. GovMut-Health đánh giá khả năng phát hiện nhóm lỗi này bằng bốn chiến lược kiểm thử đã có: kiểm thử hồi quy, kiểm thử dựa trên thuộc tính không trạng thái, kiểm thử bằng máy trạng thái và chiến lược kết hợp. Tính mới không nằm ở bản thân kiểm thử đột biến, kiểm thử dựa trên thuộc tính hay kiểm thử dựa trên mô hình, mà ở \emph{mô hình lỗi quản trị phụ thuộc lịch sử thao tác}, tập biến thể lỗi được rà soát và ánh xạ tới các bất biến của hợp đồng, cùng cách phân tích ở cấp biến thể lỗi. Đợt đánh giá đã khóa trước gồm $N=45$ biến thể lỗi không tương đương trên các tầng cổng chung và cổng cam kết; mỗi biến thể được chạy trên 16 tổ hợp phương pháp--hạt giống ngẫu nhiên, tạo 720 lượt chạy và không ghi nhận lỗi hạ tầng. Tỷ lệ phát hiện biến thể lỗi lần lượt là 35,6\% (16/45) cho kiểm thử hồi quy đơn lẻ (M0), 8,9\% (4/45) cho thuộc tính không trạng thái (M1), 13,3\% (6/45) cho kiểm thử máy trạng thái (M2) và 44,4\% (20/45) cho chiến lược kết hợp (M3). Kiểm định ghép cặp chính xác McNemar cho thấy M3 vượt trội có ý nghĩa thống kê so với M1 ($p = 3{,}05 \times 10^{-5}$) và M2 ($p = 1{,}22 \times 10^{-4}$), đồng thời bổ sung 4 biến thể bị diệt so với chỉ dùng hồi quy ($p = 0{,}125$). Phân tích 25 biến thể lỗi còn sống (tỷ lệ sống 55,6\%) làm nổi bật các điểm mù cấu trúc trong tái dựng trạng thái phía đọc và các đường dẫn cổng cam kết, chứng minh tính bổ trợ rõ nét giữa kiểm thử hồi quy và kiểm thử sinh có trạng thái, đồng thời vạch ra lộ trình thực nghiệm cho việc tổng hợp oracle trong tương lai.
\end{abstract}

\section{Giới thiệu}
Một hợp đồng quản trị có trạng thái có thể chứa lỗi chỉ lộ ra sau nhiều bước. Ví dụ, hệ thống cung cấp dữ liệu cho AI khi quyền còn hợp lệ; sau đó đồng thuận bị thu hồi; một đề xuất cũ vẫn được giữ lại; cuối cùng hệ thống thử ghi đề xuất đó vào trạng thái bền vững. Mỗi bước riêng lẻ có thể vượt qua kiểm thử, trong khi toàn bộ chuỗi vẫn vi phạm hợp đồng quản trị.

Kiểm thử dựa trên thuộc tính (property-based testing), kiểm thử dựa trên mô hình hoặc máy trạng thái (model-based/state-machine testing) và kiểm thử đột biến (mutation testing) đều là các kỹ thuật đã được nghiên cứu rộng rãi. Nghiên cứu về kiểm soát truy cập cũng đã sử dụng mô hình trạng thái và biến đổi chính sách để tạo ca lỗi. Gần đây, các hệ thống bộ nhớ cho tác tử AI và các nghiên cứu đánh giá bộ kiểm thử do LLM sinh cũng sử dụng property testing, duyệt trạng thái hữu hạn và mutation testing. Vì vậy, GovMut-Health không đề xuất một phương pháp kiểm thử cơ bản mới.

Khoảng trống nằm ở \textbf{mô hình lỗi}. Các lỗi liên quan đến độ mới của trạng thái, đồng thuận, chính sách, ràng buộc đúng chủ thể, mục đích sử dụng, ảnh chụp ngữ cảnh, cơ chế thử lại, truy vết nguồn gốc, tính nguyên tử và khả năng tái dựng kiểm toán thường phụ thuộc vào thứ tự các thao tác trước đó. GovMut-Health chủ động tạo các biến thể lỗi tại những điểm thực thi kiểm soát trong codebase, yêu cầu mỗi biến thể phải gắn với một bất biến cụ thể có thể bị phá vỡ, rồi so sánh các chiến lược kiểm thử dưới cùng một ngân sách đã khóa trước.

\section{Mô hình lỗi và đơn vị phân tích}
Tập biến thể lỗi bao phủ các nhóm như: bỏ kiểm tra phiên bản trạng thái; bỏ bước kiểm tra lại đồng thuận hoặc chính sách; cho phép ảnh chụp thuộc sai chủ thể hoặc sai mục đích; chấp nhận mã băm hoặc thời hạn không hợp lệ; tái sử dụng trạng thái dẫn xuất đã cũ; phá cơ chế chống ghi lặp; thử lại một đề xuất đã bị từ chối mà không xin lại quyền; làm mất thông tin nguồn gốc; hoặc phá tính nguyên tử và khả năng tái dựng dấu vết kiểm toán.

Một biến thể lỗi chỉ được đưa vào mẫu phân tích khi có thể chạy được và đã được rà soát là \emph{không tương đương} với chương trình gốc. Ở lần đóng băng cuối, việc sàng lọc này dùng quy trình mù với hai mô hình khóa, Gemini 3.6 Flash High và Claude Sonnet 4.6, theo cùng một tiêu chí đã cố định trước. Chỉ những biến thể mà cả hai cùng xác nhận là không tương đương (kể cả sau bước hòa giải đã quy định trước) mới được đưa vào mẫu số. Đây là rà soát bằng hai mô hình nhằm tăng khả năng tái lập, \textbf{không phải thẩm định độc lập của chuyên gia con người}. Nếu một biến thể không thể thay đổi hành vi quan sát được, hoặc không thể biên dịch hoặc chạy ổn định, lý do loại phải được ghi lại. Chỉ số chính là \emph{tỷ lệ phát hiện biến thể lỗi} (mutation score):
\[
MS(T)=\frac{K_T}{N_{exec}-N_{eq}},
\]
trong đó $K_T$ là số biến thể lỗi bị chiến lược $T$ phát hiện. Đơn vị suy luận thống kê là \textbf{một biến thể lỗi}, không phải số chuỗi thực thi hay số hạt giống ngẫu nhiên.

\section{Bốn chiến lược kiểm thử}
\textbf{M0 -- Kiểm thử hồi quy (Regression):} sử dụng bộ kiểm thử hiện có, chủ yếu bảo vệ các hành vi đã biết và các lỗi từng được quan sát.

\textbf{M1 -- Kiểm thử thuộc tính không trạng thái (Stateless properties):} sinh nhiều đầu vào độc lập và kiểm tra bất biến cục bộ, nhưng không duy trì lịch sử trạng thái qua một chuỗi nhiều bước.

\textbf{M2 -- Kiểm thử bằng máy trạng thái (State machine):} sinh các chuỗi chuyển trạng thái hợp lệ và không hợp lệ từ một mô hình trạng thái, rồi kiểm tra bất biến sau từng bước.

\textbf{M3 -- Kết hợp (Combined):} sử dụng đồng thời kiểm thử hồi quy, kiểm thử thuộc tính không trạng thái và kiểm thử bằng máy trạng thái trong cùng ngân sách đã khóa trước.

Hạt giống ngẫu nhiên, số ví dụ sinh ra, độ dài chuỗi trạng thái, phiên bản khung kiểm thử, môi trường chạy và danh sách biến thể lỗi đều được khóa trước khi đánh giá. Việc chạy nhiều hạt giống giúp kiểm tra độ ổn định và khả năng tái lập, nhưng không làm tăng số đơn vị khoa học độc lập.

\section{Kết quả}
Đợt chạy niêm phong \texttt{govmut-soict-2026-final-v2}, dựa trên phiên bản mã nguồn \texttt{ab877e04}, được thực thi trên toàn bộ 45 biến thể lỗi qua 16 cấu hình phương pháp--hạt giống (720 lượt chạy) với 0 lỗi hạ tầng.

\begin{center}
\small
\begin{tabular}{lrrrr}
\toprule
Chiến lược & Bị diệt / 45 & Tỷ lệ phát hiện & McNemar $p$ (vs.\ M3) & McNemar $p$ (vs.\ M0) \\
\midrule
M0 Hồi quy & 16 / 45 & 35,6\% & $p = 0{,}125$ & --- \\
M1 Thuộc tính không trạng thái & 4 / 45 & 8,9\% & $p = 3{,}05 \times 10^{-5}$ & $p = 0{,}00183$ \\
M2 Máy trạng thái & 6 / 45 & 13,3\% & $p = 1{,}22 \times 10^{-4}$ & $p = 0{,}0213$ \\
M3 Kết hợp & \textbf{20 / 45} & \textbf{44,4\%} & --- & $p = 0{,}125$ \\
\bottomrule
\end{tabular}
\end{center}

KTC Wilson 95\% cho tỷ lệ phát hiện lần lượt là 23,2--50,2\% (M0), 3,5--20,7\% (M1), 6,3--26,2\% (M2) và 30,9--58,8\% (M3).

M3 bao trùm toàn bộ 16 biến thể mà M0 phát hiện và thêm 4 biến thể khác ($M0\text{-only} = 0, M3\text{-only} = 4$). Bốn biến thể phát hiện thêm bởi M3 là M02-A (kiểm tra lại phiên bản chính sách trong manifest snapshot), M03-A và M03-B (kiểm tra lại phiên bản đồng thuận) và M11-A (kiểm tra liên kết phả hệ nguồn gốc trong bộ biên dịch THSS). Chênh lệch 8,9 điểm phần trăm giữa M3 và M0 chưa đạt ý nghĩa thống kê trong kiểm định ghép cặp chính xác ($p=0{,}125$). M3 phát hiện 16 biến thể mà M1 bỏ sót và 14 biến thể mà M2 bỏ sót, tạo ra các tương phản ghép cặp có ý nghĩa thống kê ($p = 3{,}05\times10^{-5}$ so với M1, $p = 1{,}22\times10^{-4}$ so với M2). Trong so sánh cặp, M0 cũng phát hiện nhiều biến thể hơn M1 ($p=0{,}00183$) và M2 ($p=0{,}0213$).

Trên toàn bộ 720 lượt chạy riêng lẻ, sổ cái cấp hạt giống ghi nhận 161 kết quả \textsc{Killed}, 559 \textsc{Survived} và 0 lỗi hạ tầng; các số đếm lượt chạy này không thay thế mẫu số cấp biến thể lỗi.

Kết quả ủng hộ tuyên bố về tính bổ trợ thay vì một thứ bậc vượt trội phổ quát giữa các phương pháp kiểm thử. M3 là phép hợp của M0, M1 và M2, do đó điểm phát hiện thô lớn hơn phản ánh độ bao phủ kết hợp trong cùng ngân sách đã khóa.

\section{Diễn giải và Đe dọa Tính Hợp lệ}
Thực nghiệm đo lường độ thỏa đáng của kiểm thử trên tập 45 biến thể đã khóa. Chiến lược kết hợp đạt tỷ lệ phát hiện biến thể cao nhất (44,4\%), nhưng lợi thế bốn biến thể so với kiểm thử hồi quy đơn lẻ chưa đạt ý nghĩa thống kê ($p = 0{,}125$). Kết luận vững chắc nhất là kiểm thử quản trị sinh tự động bổ trợ cho bộ kiểm thử hồi quy hiện có thay vì thay thế nó.

Bản thân kiểm thử hồi quy vẫn là một thành phần mạnh mẽ, tiêu diệt 16 biến thể (35,6\%) và vượt trội so với các chiến lược thuộc tính không trạng thái và máy trạng thái đơn lẻ trong tập chuẩn này. Chiến lược kết hợp bảo toàn mọi biến thể bị diệt bởi hồi quy và bổ sung thêm bốn biến thể.

Quan trọng là, 25 trong số 45 biến thể lỗi (55,6\%) vẫn sống sót qua cả bốn chiến lược. Chẩn đoán kỹ thuật mô tả phân loại các biến thể còn sống này thành bốn nhóm:
\begin{enumerate}
    \item \emph{Thiếu oracle} ($n=17$, 68,0\%), nơi các phương pháp tiếp cận được vị trí kiểm soát nhưng không dựng điều kiện vi phạm cụ thể (ví dụ: phiên bản cơ sở cũ khi áp dụng cam kết, phiên bản đồng thuận lệch trong đề xuất);
    \item \emph{Điểm mù phát lại/tái dựng} ($n=3$, 12,0\%), nơi các đường dẫn tái dựng ảnh chụp và trạng thái phía đọc chưa được kích hoạt;
    \item \emph{Thiếu đường dẫn/mục tiêu kiểm thử} ($n=3$, 12,0\%), bao gồm luồng duyệt cam kết và đường dẫn ràng buộc chéo chế độ chưa được kích hoạt; và
    \item \emph{Biến thể yếu tiềm năng} ($n=2$, 8,0\%), nơi các rào chắn bị che bởi các kiểm tra trùng lặp sâu hơn.
\end{enumerate}
Các biến thể còn sống phản ánh điểm yếu còn lại của bộ kiểm thử và cung cấp một lộ trình thực nghiệm để bổ sung các oracle tái dựng và cam kết, thay vì xem các khiếm khuyết đó là vô hại.

\section{Đe dọa tính hợp lệ}
Tính hợp lệ về cấu trúc phụ thuộc vào việc các biến thể lỗi có đại diện cho những sai sót quản trị dữ liệu sức khỏe có ý nghĩa hay không, như chấp nhận ghi dựa trên trạng thái thuốc đã cũ, bỏ qua thay đổi đồng thuận/chính sách, nhầm chủ thể hoặc mục đích, tái sử dụng ảnh chụp hết hạn/bị sửa, làm mất nguồn gốc hoặc phá tính nguyên tử giữa trạng thái và dấu vết kiểm toán. Việc xác định biến thể tương đương dùng hai mô hình khóa chứ chưa có thẩm định độc lập của chuyên gia con người, nên đây vẫn là một nguồn bất định. Tính hợp lệ nội tại chịu ảnh hưởng của tính ngẫu nhiên trong các bộ sinh ca kiểm thử; việc khóa hạt giống, ngân sách và lưu phản ví dụ đã thu gọn giúp làm phần biến thiên này minh bạch. Khả năng khái quát còn hạn chế vì GLHS là hệ thống chính được thử và tập 45 biến thể vẫn nhỏ so với các nghiên cứu mutation testing quy mô lớn. Đây là nghiên cứu về bảo đảm chất lượng phần mềm, không phải đánh giá hiệu quả lâm sàng hay tuân thủ quy định.

\section{Khả năng tái lập}
Mỗi lần chạy cuối cùng đều ghi lại mã chạy bất biến, phiên bản mã nguồn, mã băm manifest đột biến, mã băm rà soát mô hình, phiên bản khung kiểm thử, hạt giống theo thứ tự, ngân sách, siêu dữ liệu máy chủ, vết phản ví dụ chính xác, thời gian chạy thô, ma trận biến thể bị diệt và danh mục SHA-256 hiện vật. Toàn bộ ma trận chi tiết giữa từng biến thể lỗi và từng chiến lược được lưu trong \path{research/assurance_soict/results/final-analysis.json}.

\section{Mối quan hệ với công trình GLHS trước đây}
Nghiên cứu kiến trúc GLHS khảo sát khả năng duy trì liên kết ngữ cảnh sức khỏe dài hạn của mô hình với việc chấp thuận ghi bền vững. GovMut-Health đặt câu hỏi kỹ thuật phần mềm riêng biệt: chiến lược kiểm thử nào phát hiện được các khiếm khuyết quản trị phụ thuộc chuỗi thao tác được cấy vào một hệ thống như vậy? Bộ chuẩn này không tái sử dụng nhóm 64 đối tượng của mô hình GLHS làm mẫu khoa học.

\section{Kết luận}
GovMut-Health định nghĩa bộ chuẩn kiểm thử đột biến có thể thực thi cho các lỗi quản trị phụ thuộc chuỗi thao tác trong các hệ thống AI y tế có trạng thái. Đóng góp cốt lõi là phân loại lỗi đặc thù y tế, ánh xạ bất biến và so sánh ở cấp biến thể lỗi giữa các chiến lược kiểm thử đã thiết lập. Trong thực nghiệm niêm phong ($N=45$, 720 lượt chạy), chiến lược kết hợp đạt tỷ lệ phát hiện cao nhất (44,4\%, 20/45 biến thể bị diệt); mức tăng 4 biến thể so với kiểm thử hồi quy chưa có ý nghĩa thống kê ($p = 0{,}125$), trong khi mức tăng so với thuộc tính không trạng thái ($p = 3{,}05 \times 10^{-5}$) và máy trạng thái ($p = 1{,}22 \times 10^{-4}$) đều có ý nghĩa thống kê. 25 biến thể lỗi còn sống (tỷ lệ sống 55,6\%) cung cấp một lộ trình thực nghiệm để bổ sung các oracle kiểm thử tái dựng phía đọc và luồng cổng cam kết. Các kết quả này ủng hộ một chiến lược kiểm thử bổ trợ cho tập chuẩn đã khóa, thay vì tuyên bố kiểm chứng hình thức hay sự vượt trội tuyệt đối.
\end{document}
